\subsection{分裂约化李代数的包络代数中的双序}

设 \(\mathfrak{g}\) 是 \(\mathbf{Q}\) 上的约化李代数，\(\mathfrak{h}\) 是 \(\mathfrak{g}\) 的可分裂嘉当子代数，且 \(R = R(\mathfrak{g},\mathfrak{h})\)（第 2 节，第 1 号，备注 5）。

定义 1. 若 \(\mathcal{H}\) 是 \(\mathfrak{h}\) 中的一个格，则称其为相对于 \(\mathfrak{g}\) 的许可格，只要对所有 \(\alpha \in \mathrm{R}\) 都有 \(H_{\alpha} \in \mathcal{H}\) 及 \(\alpha(\mathcal{H}) \subset \mathbf{Z}\)。

备注。1) 令 B 为 R 的一组基。格 \(\mathcal{H}\) 位于 \(\mathfrak{h}\) 中。它是许可格，当且仅当 \(H_{\alpha} \in \mathcal{H}\) 且 \(\alpha(\mathcal{H}) \subset \mathbf{Z}\) 对所有 \(\alpha \in B\) 都成立。

\begin{enumerate}
\def\labelenumi{\arabic{enumi})}
\setcounter{enumi}{1}
\item
  令 \(\mathfrak{c}\) 为 \(\mathfrak{g}\) 的中心。则格 \(\mathcal{H}\) 位于 \(\mathfrak{h}\) 中，并且是许可格，当且仅当 \(\mathrm{Q}(\mathrm{R}^{\vee}) \subset \mathcal{H} \subset \mathrm{P}(\mathrm{R}^{\vee}) \oplus \mathfrak{c}\)。此时，格 \(\mathcal{H} \cap \mathcal{D}\mathfrak{g}\) 是嘉当子代数 \(\mathfrak{h} \cap \mathcal{D}\mathfrak{g}\) 中的 \(\mathcal{D}\mathfrak{g}\) 的许可格。可能存在许可格 \(\mathcal{H}\) 满足 \(\mathcal{H} \neq (\mathcal{H} \cap \mathcal{D}\mathfrak{g}) \oplus (\mathcal{H} \cap \mathfrak{c})\)（参见第 13 节，第 1 号.IX）。
\item
  若 \(\mathfrak{g}\) 是半单的，则 \(\mathfrak{h}\) 中的许可格正是其子群 \(\mathcal{H}\)，即 \(\mathfrak{h}\) 的子群，且满足 \(Q(R^{\vee})\subset \mathcal{H}\subset P(R^{\vee})\)。
\end{enumerate}

在本号余下部分，我们固定一个分裂约化李代数 \((\mathfrak{g},\mathfrak{h})\)、\(\mathrm{R}=\mathrm{R}(\mathfrak{g},\mathfrak{h})\) 的一组基 B，并且对每个 \(\alpha\in B\) 取一对 \((x_{\alpha},y_{\alpha})\)，满足

\[
y _ {\alpha} \in \mathfrak {g} ^ {- \alpha}, x _ {\alpha} \in \mathfrak {g} ^ {\alpha}, [ y _ {\alpha}, x _ {\alpha} ] = H _ {\alpha}.\tag{20}
\]

若记 \(n_{+}\)（分别为 \(n_{-}\)）为由 \(x_{\alpha}\)（分别为 \(y_{\alpha}\)）生成的 g 的子代数，则我们知道（第 3 节，第 3 号，命题 9 (iii)）

\[
\mathfrak {g} = \mathfrak {n} _ {-} \oplus \mathfrak {h} \oplus \mathfrak {n} _ {+}\tag{21}
\]

\[
\mathrm{U} (\mathfrak {g}) = \mathrm{U} (\mathfrak {n} _ {-}) \otimes_ {\mathbf {Q}} \mathrm{U} (\mathfrak {h}) \otimes_ {\mathbf {Q}} \mathrm{U} (\mathfrak {n} _ {+})\tag{22}
\]

（其中 \(\mathrm{U}(\mathfrak{g}),\ldots\) 分别是 \(\mathfrak{g},\ldots\) 的包络代数）。

记 \(\mathcal{U}_{+}\) 为一个 \(\mathbf{Z}\)-子代数，包含于 \(\mathrm{U}(\mathfrak{n}_{+})\)，并由 \(x_{\alpha}^{(n)}\)（其中 \(\alpha \in B\) 且 \(n \in \mathbf{N}\)）生成。令 W 为 R 的 Weyl 群，\(R_{+}\) 为相对于 B 的正根集合。

引理 5. (i) \(\mathcal{U}_{+}\) 是向量空间 \(\mathrm{U}(\mathfrak{n}_{+})\) 中的一个格。

\begin{enumerate}
\def\labelenumi{(\roman{enumi})}
\setcounter{enumi}{1}
\tightlist
\item
  对所有 \(\alpha \in \mathrm{B}\)，有 \(\mathcal{U}_{+} \cap \mathrm{U}(\mathfrak{g}^{\alpha}) = \bigoplus_{n \in \mathbf{N}} \mathbf{Z}x_{\alpha}^{(n)}\)。
\end{enumerate}

按定义，\(\mathcal{U}_{+}\) 作为 \(\mathbf{Z}\)-模由以下元素生成

\[
x _ {\varphi} ^ {(\mathbf {n})} = \prod_ {1 \leq i \leq r} x _ {\varphi (i)} ^ {(n (i))}
\]

其中 \(r \in \mathbf{N}\)、\(\varphi = (\varphi(i)) \in \mathrm{B}^{r}\)，且 \(\mathbf{n} = (n(i)) \in \mathbf{N}^{r}\)。给代数 \(\mathrm{U}(\mathfrak{n}_{+})\) 赋予 \(\mathrm{Q}(\mathrm{R})\) 型的分次，使每个 \(\mathfrak{g}^{\alpha} (\alpha \in \mathrm{R}_{+})\) 都是次数为 \(\alpha\) 的齐次空间。前述类型的单项式 \(x_{\varphi}^{(\mathbf{n})}\) 是次数为

\[
\sum_ {1 \leq i \leq r} n (i) \varphi (i) \in \mathrm{Q} (\mathrm{R}).
\]

具有给定次数 q 的这类单项式数目有限，并且它们在 Q 上生成 \(\mathrm{U}(\mathfrak{n}_{+})\) 的次数为 q 的齐次分量。这证明 (i)。

若 \(\alpha \in \mathrm{B}\)，则 \(\mathcal{U}_{+} \cap \mathrm{U}(\mathfrak{g}^{\alpha})\) 包含于次数为 \(\alpha\) 的倍数的齐次分量之和中；因此，根据上面的结论，\(\mathcal{U}_{+} \cap \mathrm{U}(\mathfrak{g}^{\alpha})\) 由 \(x_{\varphi}^{(\mathbf{n})}\) 生成，其中 \(\sum n(i)\varphi(i) \in \mathbf{N}\alpha\)。这迫使 \(\varphi(i) = \alpha\) 对所有 \(i\) 都成立（因为 B 是 R 的一组基），所以

\[
x _ {\varphi} ^ {(\mathbf {n})} = x _ {\alpha} ^ {(n (1))} \dots x _ {\alpha} ^ {(n (r))} = \frac {(n (1) + \cdots + n (r)) !}{n (1) ! \dots n (r) !} x _ {\alpha} ^ {(n (1) + \dots + n (r))}.
\]

因此，\(\mathcal{U}_{+} \cap \mathcal{U}(\mathfrak{g}^{\alpha}) \subset \bigoplus_{n} \mathbf{Z}x_{\alpha}^{(n)}\)，故 (ii) 成立。

证毕。

在本号余下部分，若 E 和 F 是 \(U(\mathfrak{g})\) 的 Z-子模，我们用 E.F 表示 \(U(\mathfrak{g})\) 的 Z-子模；它由乘积 ab 生成，其中 \(a \in E, b \in F\)。

命题 3. 设 \(\mathcal{H}\) 是 \(\mathfrak{h}\) 中的一个许可格。令 \(\mathcal{U}_{+},\mathcal{U}_{-},\mathcal{U}_{0}\) 为 \(\mathbf{Z}\)-子代数，包含于 \(\mathrm{U}(\mathfrak{g})\)，它们分别由元素 \(x_{\alpha}^{(n)}\)（\(\alpha \in \mathrm{B}, n \in \mathbf{N}\)）、\(y_{\alpha}^{(n)}\)（\(\alpha \in \mathrm{B}, n \in \mathbf{N}\)）、\(\binom{h}{n}\)（\(h \in \mathcal{H}, n \in \mathbf{N}\)）生成。令 \(\mathcal{U}\) 为 \(\mathbf{Z}\)-子代数 \(\mathrm{U}(\mathfrak{g})\)，由 \(\mathcal{U}_{+},\mathcal{U}_{-},\mathcal{U}_{0}\) 生成。

\begin{enumerate}
\def\labelenumi{(\roman{enumi})}
\item
  \(\mathcal{U}\) 是双代数 \(\mathrm{U}(\mathfrak{g})\) 中的双序。
\item
  有 \(\mathcal{U} = \mathcal{U}_{-}. \mathcal{U}_{0}. \mathcal{U}_{+}\)、\(\mathcal{U} \cap \mathfrak{h} = \mathcal{H}\)，并且对所有 \(\alpha \in B\)，
\end{enumerate}

\[
\mathcal {U} \cap \mathfrak {g} ^ {\alpha} = \mathbf {Z} x _ {\alpha}, \quad \mathcal {U} \cap \mathfrak {g} ^ {- \alpha} = \mathbf {Z} y _ {\alpha}.
\]

由引理 5 和命题 2，\(\mathcal{U}_{+},\mathcal{U}_{-},\mathcal{U}_{0}\) 分别是 \(\mathbf{Q}\)-代数 \(\mathrm{U}(\mathfrak{n}_{+}),\mathrm{U}(\mathfrak{n}_{-}),\mathrm{U}(\mathfrak{h})\) 中的序，并且

\[
\binom{\pm h + q}{p} \in \mathcal {U} _ {0} \quad \text { for } h \in \mathcal {H}, q \in \mathbf {Z}, p \in \mathbf {N}.\tag{23}
\]

置 \(\mathcal{L} = \mathcal{U}_{-}. \mathcal{U}_{0}. \mathcal{U}_{+} \subset \mathrm{U}(\mathfrak{g})\)。根据 (22)，\(\mathcal{L}\) 是 \(\mathrm{U}(\mathfrak{g})\) 中的一个格。由构造，

\[
\mathcal {U} _ {-}. \mathcal {L} \subset \mathcal {L}\tag{24}
\]

\[
\mathcal {L}. \mathcal {U} _ {+} \subset \mathcal {L}\tag{25}
\]

而引理 3 及 (23) 蕴含

\[
\mathcal {U} _ {0}. \mathcal {L} \subset \mathcal {L}\tag{26}
\]

\[
\mathcal {L}. \mathcal {U} _ {0} \subset \mathcal {L}.\tag{27}
\]

\[
\text {   Let   } \alpha \in \mathrm{B}, n \in \mathbf {N}, r \in \mathbf {N}, \varphi = (\varphi (i)) \in \mathrm{B} ^ {r}, \text {   and   }
\]

\[
(m (1), \dots , m (r)) \in \mathbf {N} ^ {r}.
\]

我们证明

\[
x _ {\alpha} ^ {(n)} y _ {\varphi (1)} ^ {(m (1))} \dots y _ {\varphi (r)} ^ {(m (r))} \in \mathcal {L}\tag{28}
\]

或者，鉴于 (25)，等价地证明

\[
[ x _ {\alpha} ^ {(n)}, y _ {\varphi (1)} ^ {(m (1))} \dots y _ {\varphi (r)} ^ {(m (r))} ] \in \mathcal {L}.\tag{29}
\]

我们对 \(r\) 归纳。要研究的括号是以下各项之和

\[
y _ {\varphi (1)} ^ {(m (1))} \dots y _ {\varphi (k)} ^ {(m (k))} [ x _ {\alpha} ^ {(n)}, y _ {\varphi (k + 1)} ^ {(m (k + 1))} ] y _ {\varphi (k + 2)} ^ {(m (k + 2))} \dots y _ {\varphi (r)} ^ {(m (r))}.\tag{30}
\]

当 \(\alpha \neq \varphi (k + 1)\) 时，\(x_{\alpha}\) 与 \(y_{\varphi (k + 1)}\) 对易，所以 \([x_{\alpha}^{(n)},y_{\varphi (k + 1)}^{(m(k + 1))}] = 0\)。若 \(\alpha = \varphi (k + 1)\)，则由 (17)，表达式 (30) 是以下形式的表达式之和

\[
y _ {\varphi (1)} ^ {(m (1))} \dots y _ {\varphi (k)} ^ {(m (k))} y _ {\varphi (k + 1)} ^ {(m (k + 1) - p)} \binom{q - h}{p} x _ {\alpha} ^ {(n - p)} y _ {\varphi (k + 2)} ^ {(m (k + 2))} \dots y _ {\varphi (r)} ^ {(m (r))}\tag{31}
\]

其中 \(q \in \mathbf{Z}, p \in \mathbf{N} - \{0\}, h \in \mathcal{H}\)。归纳假设连同 (24) 和 (26) 证明表达式 (31) 属于 \(\mathcal{L}\)。于是 (28) 得证。

由 (28)，\(x_{\alpha}^{(n)}\mathcal{U}_{-}\subset \mathcal{L}\)；因而根据 (25) 和 (27)，\(x_{\alpha}^{(n)}\mathcal{L}\subset \mathcal{L}\)，所以 \(\mathcal{U}_{+}.\mathcal{L}\subset \mathcal{L}\)，并且

\[
\mathcal {L}. \mathcal {L} \subset \mathcal {U} _ {-}. \mathcal {U} _ {0}. \mathcal {L} \subset \mathcal {U} _ {-}. \mathcal {L} \subset \mathcal {L}.
\]

因此，\(\mathcal{L}\) 是一个 \(\mathbf{Z}\)-子代数，包含于 \(\mathrm{U}(\mathfrak{g})\)，从而 \(\mathcal{U} = \mathcal{L}\)。若 \(c\) 是 \(\mathrm{U}(\mathfrak{g})\) 的余积，则 \(c(\mathcal{U}) \subset \mathcal{U} \otimes_{\mathbf{Z}} \mathcal{U}\)（第 2 号，命题 1）。令 \(\gamma\) 为 \(\mathrm{U}(\mathfrak{g})\) 的余单位。由于 \(\gamma(x_{\alpha}^{(n)}) = \gamma(y_{\alpha}^{(n)}) = \gamma\left(\binom{h}{n}\right) = 0\) 对 \(n > 0\) 成立，有 \(\gamma(\mathcal{U}) \subset \mathbf{Z}\)。这证明 (i)。另一方面，

\[
\mathcal {U} \cap \mathfrak {h} = \mathcal {L} \cap \mathfrak {h} = \mathcal {U} _ {0} \cap \mathfrak {h} = \mathcal {H}
\]

根据第 4 号的命题 2；类似地，

\[
\mathcal {U} \cap \mathfrak {g} ^ {\alpha} = \mathcal {U} _ {+} \cap \mathfrak {g} ^ {\alpha} = \mathbf {Z} x _ {\alpha}
\]

根据引理 5。这证明 (ii)。

备注 4. 根据第 4 节第 4 号的命题 5，存在唯一的自同构 \(\theta\)，定义在 \(\mathfrak{g}\) 上，并满足 \(\theta(x_{\alpha}) = y_{\alpha}\) 和 \(\theta(y_{\alpha}) = x_{\alpha}\) 对所有 \(\alpha \in \mathrm{B}\) 成立，且 \(\theta(h) = -h\) 对所有 \(h \in \mathfrak{h}\) 成立；有 \(\theta^2 = 1\)。由 \(\mathcal{U}\) 的构造可见，在 \(\mathrm{U}(\mathfrak{g})\) 上延拓 \(\theta\) 所得的自同构保持 \(\mathcal{U}\) 稳定。

推论 1. 置 \(\mathcal{G} = \mathcal{U} \cap \mathfrak{g}\)。则 \(\mathcal{G}\) 是李代数 \(\mathfrak{g}\) 中的一个序，并且在 \(\theta\) 下稳定。有 \(\mathcal{G} = \mathcal{H} + \sum_{\alpha \in \mathrm{R}} (\mathcal{G} \cap \mathfrak{g}^{\alpha})\)。对于所有 \(\alpha \in \mathrm{B}\) 及所有 \(n \in \mathbf{N}\)，映射 \((\operatorname{ad} x_{\alpha})^{n} / n!\)、\((\operatorname{ad} y_{\alpha})^{n} / n!\) 都保持 \(\mathcal{U}\) 和 \(\mathcal{G}\) 稳定。

第一个断言显然。第二个断言由考察类型为 \(Q(R)\) 的分次（定义在 \(U(\mathfrak{g})\) 和 U 上）得到。第三个断言由第 5 号的引理 2 得到。

推论 2. 令 \(w \in \mathrm{W}\)。存在初等自同构 \(\varphi\)，它作用于 \(\mathfrak{g}\)，与 \(\theta\) 对易，保持 \(\mathcal{G}\) 和 \(\mathcal{U}\) 稳定，并延拓 \(w\)。

只需处理 \(w\) 形如 \(s_{\alpha}\)（\(\alpha \in \mathrm{B}\)）的情形。首先注意，\(\operatorname{ad} x_{\alpha}\) 和 \(\operatorname{ad} y_{\alpha}\) 在 \(\mathrm{U}(\mathfrak{g})\) 上局部幂零；换言之，对每个 \(u \in \mathrm{U}(\mathfrak{g})\)，存在一个整数 \(n\)，使得 \((\operatorname{ad} x_{\alpha})^n u = (\operatorname{ad} y_{\alpha})^n u = 0\)。因此可以定义自同构 \(e^{\operatorname{ad} x_{\alpha}} = \sum_{n=0}^{\infty} \frac{1}{n!} (\operatorname{ad} x_{\alpha})^n\) 和 \(e^{\mathrm{ad}y_{\alpha}}\) 作用于 \(\mathrm{U}(\mathfrak{g})\)；立即验证，这些 \(\mathrm{U}(\mathfrak{g})\) 的自同构保持 \(\mathcal{U}\) 稳定。置 \(\varphi_1 = e^{\mathrm{ad}x_\alpha}e^{\mathrm{ad}y_\alpha}e^{\mathrm{ad}x_\alpha}\)、\(\varphi_2 = e^{\mathrm{ad}y_\alpha}e^{\mathrm{ad}x_\alpha}e^{\mathrm{ad}y_\alpha}\)。有 \(\varphi_1|\mathfrak{g} = \varphi_2|\mathfrak{g}\)（第 2 节，第 2 号，公式 (1)），故 \(\varphi_1 = \varphi_2\)。置 \(\varphi_1 = \varphi_2 = \varphi\)。有 \(\theta \varphi \theta^{-1} = \varphi\)，所以 \(\theta\) 与 \(\varphi\) 对易。另一方面，根据第 2 节第 2 号引理 1，有 \(\varphi |\mathfrak{h} = w\)。

推论 3. 令 \(\alpha \in \mathrm{R}\)。若 \(x \in \mathcal{G} \cap \mathfrak{g}^{\alpha}\) 且 \(n \in \mathbf{N}\)，则 \(x^{(n)} \in \mathcal{U}\)，并且 \((\operatorname{ad} x)^n / n!\) 保持 \(\mathcal{G}\) 和 \(\mathcal{U}\) 稳定。

当 \(\alpha \in \mathrm{B}\) 时，由 \(\mathcal{U}\) 的构造及推论 1，这显然成立。在一般情形，存在 \(w \in \mathrm{W}\) 使 \(w(\alpha) \in \mathrm{B}\)（第 VI 章，第 1 节，第 5 号，命题 15）。根据推论 2，存在自同构 \(\varphi\)，作用于 \(\mathfrak{g}\)，保持 \(\mathcal{G}\) 和 \(\mathcal{U}\) 稳定，并将 \(\mathfrak{g}^{\alpha}\) 送到 \(\mathfrak{g}^{w(\alpha)}\)，故由结构迁移得到该推论。

推论 4. 存在一个谢瓦莱系 \((X_{\alpha})_{\alpha \in \mathrm{R}}\)，它位于 \((\mathfrak{g},\mathfrak{h})\) 中（第 2 节，第 4 号，定义 3），并满足 \(X_{\alpha} = x_{\alpha}\) 且 \(X_{-\alpha} = y_{\alpha}\) 对 \(\alpha \in \mathrm{B}\) 成立。对于具有这些性质的任意谢瓦莱系 \((X_{\alpha}^{\prime})_{\alpha \in \mathrm{R}}\)，以及所有 \(\alpha \in \mathrm{R}\)，\(X_{\alpha}^{\prime}\) 都是 \(\mathcal{G} \cap \mathfrak{g}^{\alpha}\) 的一组基。

当 \(\alpha \in \mathrm{B}\) 时，置 \(X_{\alpha} = x_{\alpha}, X_{-\alpha} = y_{\alpha}\)。当 \(\alpha \in \mathrm{R}_{+} - \mathrm{B}\) 时，选取 \(w \in \mathrm{W}\) 使 \(w(\alpha) \in \mathrm{B}\)，并选取自同构 \(\varphi\)，它作用于 \(\mathfrak{g}\)，使得 \(\theta \varphi = \varphi \theta\)、\(\varphi(\mathcal{G}) = \mathcal{G}\)，且 \(\varphi(h) = w^{-1}(h)\) 对 \(h \in \mathfrak{h}\) 成立（推论 2）；置 \(X_{\alpha} = \varphi(x_{w(\alpha)})\)、\(X_{-\alpha} = \varphi(y_{w(\alpha)})\)。则

\[
[ X _ {- \alpha}, X _ {\alpha} ] = \varphi ([ y _ {w (\alpha)}, x _ {w (\alpha)} ]) = \varphi (H _ {w (\alpha)}) = w ^ {- 1} (H _ {w (\alpha)}) = H _ {\alpha}
\]

\[
\theta (X _ {\alpha}) = \theta \varphi (x _ {w (\alpha)}) = \varphi \theta (x _ {w (\alpha)}) = \varphi (y _ {w (\alpha)}) = X _ {- \alpha}
\]

所以 \((X_{\alpha})_{\alpha \in \mathrm{R}}\) 是一个谢瓦莱系。此外，

\[
\mathcal {G} \cap \mathfrak {g} ^ {\alpha} = \varphi (\mathcal {G} \cap \mathfrak {g} ^ {w (\alpha)}) = \varphi (\mathbf {Z} x _ {w (\alpha)}) = \mathbf {Z} X _ {\alpha}\tag{32}
\]

\[
\mathcal {G} \cap \mathfrak {g} ^ {- \alpha} = \varphi (\mathcal {G} \cap \mathfrak {g} ^ {- w (\alpha)}) = \varphi (\mathbf {Z} y _ {w (\alpha)}) = \mathbf {Z} X _ {- \alpha}.\tag{33}
\]

令 \((X'_{\alpha})_{\alpha \in \mathrm{R}}\) 为一个谢瓦莱系，满足 \(X'_{\alpha} = x_{\alpha}, X'_{-\alpha} = y_{\alpha}\) 对 \(\alpha \in \mathrm{B}\) 成立。令 \(\mathrm{S}\) 为满足以下条件的 \(\alpha \in \mathrm{R}\) 所成的集合：\(X'_{\alpha} = \pm X_{\alpha}\)。根据第 2 节第 4 号命题 7，\(\mathrm{S}\) 是根的闭集。由于 \(\mathrm{S} \supset \mathrm{B} \cup (-\mathrm{B})\)，有 \(\mathrm{S} = \mathrm{R}\)（第 VI 章，第 1 节，第 6 号，命题 19）。因此，根据 (32) 和 (33)，有 \(\mathcal{G} \cap \mathfrak{g}^{\alpha} = \mathbf{Z}X'_{\alpha}\) 对所有 \(\alpha \in \mathrm{R}\) 成立。

备注。5) 令 \((X_{\alpha})_{\alpha \in \mathrm{R}}\) 为上面构造的谢瓦莱系。若 \(\alpha, \beta, \alpha + \beta \in \mathrm{R}\)，并置 \([X_{\alpha}, X_{\beta}] = \mathrm{N}_{\alpha, \beta} X_{\alpha + \beta}\)，则有 \([X_{\alpha}, X_{\beta}] \in \mathcal{G} \cap \mathfrak{g}^{\alpha + \beta}\)，从而重新得到 \(\mathrm{N}_{\alpha, \beta} \in \mathbf{Z}\) 这一事实（参见第 2 节，第 4 号，命题 7）。

\begin{enumerate}
\def\labelenumi{\arabic{enumi})}
\setcounter{enumi}{5}
\tightlist
\item
  我们在此过程中还得到谢瓦莱系存在性的新证明（参见第 4 节，第 4 号，命题 5 的推论），且该证明独立于第 2 节的引理 4。
\end{enumerate}

